% 第2章 基本数据类型（教材第2章 2.1--2.8）
\lstset{language=C}

\section{常量与变量}

C 语言程序处理的数据有\textbf{常量}和\textbf{变量}两种形式。常量在程序运行过程中值不变，变量则对应内存中一个可读写的存储单元。

\subsection{常量}

常量按其表示的数据种类分为四类：整型常量、实型常量、字符型常量、字符串字面量。

\subsubsection{整型常量}

整型常量可以用四种进制书写，区别在于前缀：

\begin{table}[htbp]
\centering
\caption{整型常量的进制表示}
\small
\begin{tabularx}{\textwidth}{@{}llXl@{}}
\toprule
进制 & 前缀 & 示例 & 十进制值 \\
\midrule
十进制 & 无 & \texttt{123}、\texttt{-45} & 123、$-45$ \\
八进制 & 数字 \texttt{0} & \texttt{017} & 15 \\
十六进制 & \texttt{0x} 或 \texttt{0X} & \texttt{0x1F} & 31 \\
二进制 & 无（C99 起可用 \texttt{0b}，教材未采用） & —— & —— \\
\bottomrule
\end{tabularx}
\end{table}

书写整型常量时要注意后缀表示类型：\texttt{123L} 是长整型，\texttt{123U} 是无符号整型，\texttt{123UL} 是无符号长整型。八进制常量以 \texttt{0} 开头这一点极易出错：\texttt{017} 不是十七而是十五，\texttt{0} 后面若误写了 \texttt{8} 或 \texttt{9}（如 \texttt{018}）则是非法常量。

\subsubsection{实型常量}

实型常量有\textbf{小数形式}和\textbf{指数形式}两种写法：

\begin{itemize}
  \item 小数形式：必须带小数点，如 \texttt{3.14}、\texttt{.5}、\texttt{5.}、\texttt{-0.001}；
  \item 指数形式：由尾数、字母 \texttt{E}（或 \texttt{e}）和指数三部分组成，如 \texttt{1.23456e4} 表示 $1.23456\times10^{4}$，\texttt{2.5E-3} 表示 $2.5\times10^{-3}$。
\end{itemize}

指数形式有两条硬规定：\texttt{E} 前后的\textbf{尾数和小数部分都不能省略}，即 \texttt{e4}、\texttt{1.2e} 都是非法的；指数必须是整数，\texttt{1.2e1.5} 也是非法的。实型常量默认为 \texttt{double} 型，加后缀 \texttt{f}（如 \texttt{3.14f}）则为 \texttt{float} 型。

\subsubsection{字符常量与字符串字面量}

字符常量是用一对\textbf{单引号}括起来的单个字符，如 \texttt{'a'}、\texttt{'A'}、\texttt{'5'}。字符在内存中存的是它的 ASCII 码值，因此 \texttt{'a'} 与整数 97 在数值上等价，可以参与算术运算。

反斜杠开头的\textbf{转义字符}用来表示无法直接写出的控制字符或特殊符号：

\begin{table}[htbp]
\centering
\caption{常用转义字符}
\small
\begin{tabularx}{\textwidth}{@{}llX@{}}
\toprule
转义字符 & 含义 & 说明 \\
\midrule
\texttt{\textbackslash n} & 换行 & 光标移到下一行行首 \\
\texttt{\textbackslash t} & 水平制表 & 跳到下一个制表位 \\
\texttt{\textbackslash r} & 回车 & 光标回到本行行首 \\
\texttt{\textbackslash b} & 退格 & 光标左移一格 \\
\texttt{\textbackslash\textbackslash} & 反斜杠 & 输出一个 \texttt{\textbackslash} \\
\texttt{\textbackslash'} / \texttt{\textbackslash"} & 单/双引号 & 在字符串中输出引号本身 \\
\texttt{\textbackslash 0} 或 \texttt{\textbackslash 000} & 空字符 & ASCII 值为 0，字符串结束标志 \\
\bottomrule
\end{tabularx}
\end{table}

字符串字面量（常量）是用一对\textbf{双引号}括起来的若干字符，如 \texttt{"Hello"}。两者的关键差别：\texttt{'a'} 是字符常量，占 1 个字节；\texttt{"a"} 是字符串，占 2 个字节（字符 \texttt{a} 加一个 \texttt{'\textbackslash 0'} 结束标志）。所以 \texttt{'a'} 与 \texttt{"a"} 不能混用，单引号只能括一个字符，双引号不能用来表示字符。

\subsection{变量}

变量在使用之前\textbf{必须先定义}，定义变量的一般形式为
\[
\text{类型关键字}\quad\text{变量名};
\]
变量名是用户定义的标识符，用于标识内存中一个具体的存储单元。命名必须遵守：

\begin{enumerate}
  \item 第一个字符必须是字母（A--Z、a--z）或下划线 \texttt{\_}；
  \item 后续字符可以是字母、数字或下划线；
  \item 不能包含空格或其他特殊字符；
  \item 不能使用 C 的关键字（如 \texttt{int}、\texttt{return}、\texttt{if}）；
  \item \textbf{区分大小写}：\texttt{newValue} 与 \texttt{newvalue} 是两个不同的变量。
\end{enumerate}

C 语言允许定义变量的同时\textbf{初始化}（赋初值）；一条语句可同时定义多个同类型变量，之间用逗号分隔。但要注意下面两种写法的区别：

\begin{lstlisting}
int a = 0, b = 0, c = 0;   /* 正确：三个变量各自初始化 */
int a = b = c = 0;         /* 错误：b、c 此时尚未定义 */
\end{lstlisting}

还有一个易被忽视的约束：\textbf{变量必须在可执行语句之前定义}。把 \texttt{int n;} 写在 \texttt{printf} 之后，C89 编译器会报错（C99 之后允许，但为了可移植性，仍应统一写在函数开头）。

\corepoint{定义变量时问两个问题：它要存什么类型的数据（决定类型关键字）、它的值从哪里来（决定是否需要初始化）。未初始化的局部变量里是\textbf{残留的垃圾值}，不是 0——这是“结果莫名其妙”的最常见来源。}

\section{简单的屏幕输出}

一个能运行的最小 C 程序结构如下：

\begin{lstlisting}
#include <stdio.h>
int main(void)
{
    printf("Hello, world!\n");
    return 0;
}
\end{lstlisting}

第一行以 \texttt{\#} 开头且不以分号结尾，它不是 C 语句，而是\textbf{编译预处理命令}。尖括号内的是\textbf{头文件}，\texttt{stdio.h} 中声明了标准输入/输出函数的原型；不包含它而直接调用 \texttt{printf()}，编译器会警告甚至报错。

关于 \texttt{printf()} 需要先建立三点认识：

\begin{itemize}
  \item 它的作用是\textbf{按格式}把数据输出到屏幕，最简单的用法是输出一个字符串；
  \item 用一对双引号括起来的若干字符称为字符串，字符串中可以包含转义字符；
  \item \texttt{main} 函数的返回值 \texttt{0} 表示程序正常结束，这个 0 会返回给操作系统。
\end{itemize}

\section{数据类型}

C 语言把数据类型分为基本类型、构造类型、指针类型和空类型四大类。之所以要区分类型，是因为不同类型的数据\textbf{取值范围不同、表示形式不同、可参与的运算也不同}。

\begin{table}[htbp]
\centering
\caption{C 语言的数据类型分类}
\footnotesize
\setlength{\tabcolsep}{4pt}
\begin{tabularx}{\textwidth}{@{}p{1.7cm}p{2.4cm}Xp{1.5cm}@{}}
\toprule
分类 & 关键字 & 变量声明实例 & 详述章节 \\
\midrule
\multirow{10}{*}{基本类型}
 & 基本整型 \texttt{int} & \texttt{int a;} & 第2章 \\
 & 长整型 \texttt{long int} & \texttt{long int a;} 或 \texttt{long a;} & 第4章 \\
 & 长长整型 \texttt{long long} & \texttt{long long int a;} 或 \texttt{long long a;} & 第4章 \\
 & 短整型 \texttt{short int} & \texttt{short int a;} 或 \texttt{short a;} & 第4章 \\
 & 无符号整型 \texttt{unsigned} & \texttt{unsigned int a;} 等 & 第4章 \\
\cmidrule(l){2-4}
 & 单精度实型 \texttt{float} & \texttt{float a;} & 第4章 \\
 & 双精度实型 \texttt{double} & \texttt{double a;} & 第4章 \\
 & 长双精度实型 \texttt{long double} & \texttt{long double a;} & 第4章 \\
\cmidrule(l){2-4}
 & 字符型 \texttt{char} & \texttt{char a;} & 第4章、第10章 \\
 & 枚举类型 \texttt{enum} & \texttt{enum response\{no, yes\};} & 第12章 \\
\midrule
\multirow{3}{*}{构造类型}
 & 数组 & \texttt{int score[10];} & 第8章 \\
 & 结构体 \texttt{struct} & \texttt{struct date\{int y; int m; int d;\};} & 第12章 \\
 & 共用体 \texttt{union} & \texttt{union\{int i; char c;\} u;} & 第12章 \\
\midrule
指针类型 & —— & \texttt{int *ptr;}\quad \texttt{char *pStr;} & 第9--11章 \\
空类型 & \texttt{void} & \texttt{void Sort(int a[], int n);} & 第8章、第11章 \\
\bottomrule
\end{tabularx}
\end{table}

几点补充：定义整型变量时若不指定 \texttt{unsigned}，隐含就是\textbf{有符号型}，而 \texttt{signed} 通常省略不写；\texttt{char} 型在算术运算中按整型处理（存放字符的 ASCII 码）；\texttt{void} 表示“无类型”，常用作无返回值的函数返回类型或无类型的指针。

\section{如何计算变量或数据类型所占内存空间的大小}

内存空间的基本单位是\textbf{字节}（byte），1 字节 = 8 \textbf{位}（bit）。位是存储信息的最小单位，字节是编址的基本单位。

\texttt{sizeof} 是 C 语言的\textbf{关键字}，不是函数名；\texttt{sizeof()} 是专门用于计算指定数据类型字节数的\textbf{运算符}。它有两个用法：

\begin{lstlisting}
sizeof(int)      /* 计算 int 型数据所占的字节数 */
sizeof(a)        /* 计算变量 a 所占的字节数 */
\end{lstlisting}

\texttt{sizeof} 在\textbf{编译时}求值，不增加程序运行时间开销，但能显著增强程序的可移植性——它是“绝不对变量所占字节数想当然”这一规则的工具保障。

\begin{table}[htbp]
\centering
\caption{常见数据类型所占字节数（典型取值，以 \texttt{sizeof} 实测为准）}
\small
\begin{tabularx}{\textwidth}{@{}lccX@{}}
\toprule
类型 & 32 位环境字节数 & 64 位环境字节数 & 说明 \\
\midrule
\texttt{char} & 1 & 1 & 恒为 1 字节，C 标准规定 \\
\texttt{short} & 2 & 2 & 长度不长于 \texttt{int} \\
\texttt{int} & 4 & 4 & 通常与机器字长相同 \\
\texttt{long} & 4 & 4（Windows）/8（Linux） & 长度不短于 \texttt{int} \\
\texttt{long long} & 8 & 8 & C99 起 \\
\texttt{float} & 4 & 4 & 约 7 位有效数字 \\
\texttt{double} & 8 & 8 & 约 15--16 位有效数字 \\
\texttt{long double} & 12 或 16 & 16 & 因编译器而异 \\
指针 & 4 & 8 & 与地址总线宽度一致 \\
\bottomrule
\end{tabularx}
\end{table}

C 标准并没有规定各种整型数据在内存中占多少字节，只要求\textbf{长整型的长度不短于基本整型，短整型的长度不长于基本整型}。同一类型在不同编译器和不同计算机系统中所占字节数可能不同，因此跨平台程序必须用 \texttt{sizeof} 而不是硬编码的常数。

\section{变量的赋值和赋值运算符}

赋值运算符 \texttt{=} 用于给变量赋值，由赋值运算符及其两侧操作数组成的表达式称为\textbf{赋值表达式}。在赋值表达式后加分号即构成\textbf{表达式语句}：

\begin{lstlisting}
x = 1;        /* 赋值表达式 + 分号 = 赋值语句 */
\end{lstlisting}

赋值运算有严格的\textbf{方向性}：把右侧表达式（右值）的值赋给左侧的变量，并\textbf{不表示}“等号两侧的值相等”。因此等号左侧只能是标识一个特定存储单元的变量名，\texttt{1 = x}、\texttt{x + 1 = y} 都是非法的。

几条与运算规则相关的结论：

\begin{itemize}
  \item \textbf{优先级}：算术运算符的优先级高于赋值运算符，所以 \texttt{x = a + b} 先算加法再赋值；
  \item \textbf{赋值表达式的值}：被规定为运算完成后左操作数的值，类型与左操作数相同。因此 \texttt{x = x + 1} 的值就是 \texttt{x} 加 1 之后的值；
  \item \textbf{结合性}：赋值运算符是\textbf{右结合}的（自右向左），而需要两个操作数的算术运算符是左结合的。
\end{itemize}

右结合性直接解释了多重赋值的执行顺序：

\begin{lstlisting}
a = b = c = 0;          /* 等价于 a = (b = (c = 0)); */
\end{lstlisting}
执行时从右向左，把最右侧表达式的值依次赋给左侧的变量，常用于给多个变量赋相同初值。而 \texttt{x / y * z} 因为是左结合，等价于 \texttt{(x / y) * z} 而非 \texttt{x / (y * z)}——两种结合性混在一起是初学阶段最典型的错误来源。

\section{本章扩充内容}

\subsection{有符号整数和无符号整数}

两者的区别在于\textbf{怎样解释整数的最高位}。对无符号整数，最高位被解释为数据位；对有符号整数，最高位被解释为\textbf{符号位}：0 表示正数，1 表示负数。

由于符号位占去 1 位，同样字节数的有符号整数所能表示的最大值只有无符号整数的一半。以 2 字节（16 位）为例：

\begin{table}[htbp]
\centering
\caption{2 字节整型数的取值范围}
\small
\begin{tabularx}{\textwidth}{@{}lXX@{}}
\toprule
类型 & 取值范围 & 说明 \\
\midrule
有符号 \texttt{short} & $-32768 \sim 32767$ & 最小值 $-2^{15}$，最大值 $2^{15}-1$ \\
无符号 \texttt{unsigned short} & $0 \sim 65535$ & 最大值 $2^{16}-1$ \\
\bottomrule
\end{tabularx}
\end{table}

负数在计算机中以\textbf{二进制补码}（complement）形式存储。由原码求补码的规则是：\textbf{符号位不变}，其余各位取反得反码，再在末位加 1。例如 $-1$ 的原码为 \texttt{1000 0000 0000 0001}，反码为 \texttt{1111 1111 1111 1110}，加 1 后得补码 \texttt{1111 1111 1111 1111}（全 1）。正数的原码、反码、补码三者相同。

用补码表示有两条好处：

\begin{enumerate}
  \item \textbf{减法可化为加法}。计算 $7-6$ 时，把 $+7$ 的补码与 $-6$ 的补码相加，舍去最高位的进位，结果正是 $+1$，因此硬件只需一个加法器；
  \item \textbf{0 的表示唯一}。若用原码，$+0$ 与 $-0$ 的编码不同；而按补码计算，$-0$ 的反码加 1 后进位被舍去，又变回全 0，与 $+0$ 一致。
\end{enumerate}

\corepoint{全 1 的 16 位编码，按有符号数解释是 $-1$，按无符号数解释是 65535。\textbf{同一串二进制位，类型不同，含义就不同}——这正是数据类型存在的意义。把无符号数当有符号数用（或反之），是“输出莫名其妙的大数/负数”的根本原因。}

\subsection{实型数据在内存中的存储格式}

实数无论写成小数形式还是指数形式，在计算机内部都按\textbf{浮点形式}存储。浮点形式是相对于定点形式而言的：定点表示中小数点位置固定（整数是小数点固定在数值位之后的特例，表示范围很小、很死板），而浮点数的小数点位置\textbf{可以浮动}，因此能表示很大或很小的数。

浮点数把实数分成\textbf{阶码}和\textbf{尾数}两部分：
\[
N = S \times r^{J}
\]
其中 $S$ 为尾数（正负均可，一般规定为纯小数形式），$J$ 为阶码（正负均可，但必须是整数），$r$ 为基数。对二进制数而言 $r=2$，即 $N=S\times 2^{J}$。存储时依次存放阶码符号、阶码数值、尾数符号、尾数数值。

两条必须记住的结论：

\begin{itemize}
  \item \textbf{阶码所占的位数决定实数的表数范围}（能表示多大、多小的数）；
  \item \textbf{尾数所占的位数决定实数的精度}（有效数字有几位），尾数的符号决定实数的正负。
\end{itemize}

用更多的位存尾数可以提高精度但缩小表数范围；用更多的位存阶码可以扩大表数范围但降低精度。因此\textbf{浮点数并非真正意义上的实数，只是其在某种范围内的近似值}。C 标准没有明确规定三种浮点类型的阶码和尾数各占多少位，不同编译器分配不同。

\subsection{单精度实型和双精度实型的有效位数}

把同一个实型常量分别赋给 \texttt{float} 与 \texttt{double} 变量并输出，末尾几位数字会不同。原因在于二者能接收的\textbf{有效数字位数}不同：\texttt{double} 型变量一般可接收 15--16 位有效数字，而 \texttt{float} 型变量仅能接收约 7 位有效数字，有效数字之后输出的数字都是不准确的。

由此得到一条实践规则：\textbf{把 \texttt{double} 数据赋给 \texttt{float} 变量时可能发生数据截断，产生舍入误差}。需要较高精度（如累加求和、迭代运算）时应使用 \texttt{double}；只有确实要节省空间或与特定接口对接时才用 \texttt{float}。

\section{本章知识点小结}

\begin{itemize}
  \item 数据分常量与变量；常量有整型、实型、字符型、字符串字面量四类；变量必须先定义后使用，区分大小写；
  \item 整型常量四种进制靠前缀区分（\texttt{0} 八进制、\texttt{0x} 十六进制），实型常量有小数与指数两种写法；
  \item 字符常量用单引号、字符串用双引号，字符串末尾隐含 \texttt{'\textbackslash 0'}；
  \item 数据类型分基本、构造、指针、空类型；类型决定取值范围、存储形式与可参与的运算；
  \item \texttt{sizeof} 是编译时求值的运算符，用于计算类型或变量的字节数，是可移植性的保障；
  \item 赋值运算符具有方向性和右结合性，赋值表达式的值等于运算后左操作数的值；
  \item 内存模型：负数用补码存储，阶码决定表数范围，尾数决定精度，浮点数是近似值。
\end{itemize}

\section{本章常见错误小结}

\begin{table}[htbp]
\centering
\caption{第 2 章常见错误与纠正}
\small
\begin{tabularx}{\textwidth}{@{}lXl@{}}
\toprule
易错写法 & 错在哪里 & 纠正 \\
\midrule
\texttt{int newVa1ue;} 后写 \texttt{newvalue = 0;} & 变量名区分大小写，等于未定义就使用 & 名字前后一致 \\
\texttt{printf} 之后写 \texttt{int n;} & 变量必须在可执行语句之前定义 & 定义集中写在函数开头 \\
\texttt{int a = b = c = 0;} & \texttt{b}、\texttt{c} 此时尚未定义 & \texttt{int a=0,b=0,c=0;} \\
\texttt{float x = 3.5;} & 常量类型为 \texttt{double}，赋值时发生转换 & 用 \texttt{3.5f} 或改变量为 \texttt{double} \\
\texttt{'a'} 与 \texttt{"a"} 混用 & 单引号是字符，双引号是字符串 & 字符用单引号 \\
\texttt{017} 当作十七 & 前导 0 表示八进制，实为 15 & 需要十进制就不要写前导 0 \\
变量未初始化直接参与运算 & 局部变量中是残留垃圾值，不是 0 & 定义时显式初始化 \\
\bottomrule
\end{tabularx}
\end{table}
